\chapter{Сколько нужно входов}
% полная версия: chapters/19_dendrites.tex

У нейрона есть ветвистые входные отростки --- дендриты. У одних нейронов это целое дерево с тысячами контактов, у других --- одна-две ветки или вовсе ничего. Для инженера число входов --- параметр схемы, и он подобран под задачу.

\section*{От нуля до ста тысяч}

У датчиков прикосновения и боли дендритов нет совсем: провод просто тянется от кожи в спинной мозг, а датчик сидит на самом его конце. Это линия с датчиком на дальнем краю.

\pencilfig{19}{\pencilart{19}}{Много входов --- сумматор; один огромный вход --- быстрый и точный ретранслятор.}

У нейронов обоняния и у нейронов, передающих сигнал от датчиков глаза и уха, --- один дендрит. Один вход, один выход: повторитель, буфер.

У детекторов направления звука --- два дендрита, смотрящие в разные стороны, по одному на каждое ухо. Разнеся два входа по разным веткам, нейрон делает «и» строже: две порции тока в одной ветке мешали бы друг другу.

У мелких нейронов схемы обучения движениям --- около четырёх коротких дендритов, каждый хватает свой вход. Таких нейронов десятки миллиардов, и каждый складывает всего четыре сигнала. Зато вместе они раскладывают вход на миллионы комбинаций, а большой выходной нейрон с сотней тысяч входов выбирает из них нужные.

И наконец, нейроны с огромными деревьями и тысячами входов --- сумматоры и классификаторы, у которых отдельные ветки работают почти как самостоятельные маленькие схемы.

\section*{Почему так}

Всё решает зарядка. Маленький нейрон с немногими входами заряжается от одного крупного соединения за доли миллисекунды: такой нейрон точен во времени и надёжно передаёт сигнал дальше. Большое дерево одно соединение не зарядит никогда --- заряд утечёт раньше. Ему нужны десятки сигналов, пришедших вместе, зато оно может их складывать и сравнивать. Мало входов и крупные соединения --- для точного времени. Много входов --- для счёта.

\pencilfig{128}{%
% charging: a small neuron reaches threshold from one large synapse at once; a big tree barely moves
% from one input and needs many inputs arriving together
\foreach \dx/\t in {0/{мало входов}, 4.3/{большое дерево}}{
  \draw[pen] (\dx,0) -- (\dx+3.6,0);
  \draw[pcGray, densely dashed, line width=0.5pt] (\dx,0.9) -- (\dx+3.6,0.9);
  \node[font=\small\itshape, text=pcGray, anchor=south] at (\dx+1.8,1.75) {\t};}
\node[lbl, anchor=east] at (-0.05,0.9) {порог};
% small neuron
\draw[pcBlue, line width=2pt] (0.6,-0.15) -- (0.6,-0.6);
\draw[penlite=pcBlue] (0,0) -- (0.6,0) -- plot[domain=0.6:0.8, samples=12] (\x,{1.3*(1-exp(-(\x-0.6)/0.18))}) -- (0.8,1.6) -- (0.84,0) -- (3.5,0);
\node[lbl, text=pcRed, anchor=west] at (0.85,1.45) {щелчок сразу};
\node[lbl, text=pcBlue, anchor=north] at (0.6,-0.62) {один крупный вход};
% big tree
\draw[pcBlue, line width=0.6pt] (4.8,-0.15) -- (4.8,-0.45);
\foreach \s in {6.15,6.2,6.25,6.3,6.35,6.4,6.45,6.5}{\draw[pcBlue, line width=0.6pt] (\s,-0.15) -- (\s,-0.45);}
\draw[penlite=pcBlue] (4.3,0) -- (4.8,0) -- plot[domain=4.8:6.15, samples=20] (\x,{0.16*((\x-4.8)/0.15)*exp(1-(\x-4.8)/0.15)}) -- plot[domain=6.15:6.62, samples=14] (\x,{1.25*(1-exp(-(\x-6.15)/0.4))}) -- (6.62,1.6) -- (6.66,0) -- (7.9,0);
\node[lbl, anchor=south] at (4.95,0.2) {почти ничего};
\node[lbl, text=pcBlue, anchor=north] at (6.33,-0.47) {много сразу};
}{Маленький нейрон заряжается от одного крупного соединения за доли миллисекунды и щёлкает сразу. Большое дерево от одного входа почти не меняется: ему нужны десятки сигналов, пришедших вместе.}

Так мы получили третью сортировку нейронов: по веществу, по прибору и по числу входов. Каждая видит одну и ту же машину с новой стороны.

\fullversion{19}
